\section{Proofs for Cascades 3 and 4}\label{app:cascade-cuts}

This appendix proves \cref{thm:cascade-cuts} and \cref{rem:cascade-widths}. Throughout, $d=2^k\ge32$, $b\ge1$, $L=d^2b$,
$m_r=L/2^r$, $\mathrm{sh}_r=\lfloor m_r/3\rfloor+b$ and $r_{\max}=k-1$, so that $m_r\ge m_{r_{\max}}=2db\ge64$ for $r\le r_{\max}$. Let
$A=[db,3db)$ be the set of input positions of rows~$1$ and~$2$. As $A_1$ is the set of all rows, $\Qstar(\{1,2\},A_1)=A$, and
$\Lin(\{1,2\})$ is the set of steps at which the input head reads a position of $A$. The traces are those of the constructors
\texttt{build\_c3\_ten\_entry\_cascade} and \texttt{build\_c4\_no\_overwrite\_cascade} in \texttt{experiments/src/schedules.rs}; the head
motions below are read off their loops, and we also checked them on the exported traces for $d\le128$. A \emph{region} is an interval
of cells of one tape, and the \emph{offset} of a cell in a region is its position minus the first position of the region.

\subsection{Taint closures}

For a work tape $j$ and a step $t$, let $\prev_j(t)$ be the last step before $t$ at which head $j$ scans the cell it scans at step $t$,
if there is one; the cell edge into $t$ of tape $j$ is then $\prev_j(t)\to t$. The \emph{wait age} $\omega_j(t)$ is $0$ if $t=0$ or
head $j$ scans different cells at steps $t-1$ and $t$, and $\omega_j(t-1)+1$ otherwise.

\begin{definition}[Taint closure]\label{def:taint}
Fix a set $C_0$ of cell edges and an order $h\in\{1,2\}$. A step $t$ is \emph{tainted} if it reads a position of $A$, or if for some
tape $j$ the cell edge $p\to t$ with $p=\prev_j(t)$ is not in $C_0$, the step $p$ is tainted, and either $p\le t-2$ or $\omega_j(t)\le h$.
A \emph{run} is a maximal interval of consecutive tainted steps. The \emph{cut} $C$ consists of $C_0$ and of the state edge and the
waits that leave the last step of each run. A cell of tape $j$ is \emph{live tainted} at step $t$ if its last visit $p<t$ is tainted and
the edge $p\to\mathrm{next}_j(p)$ is not in $C_0$.
\end{definition}

A step $t$ with $p=\prev_j(t)=t-1$ is a \emph{short wait} if $\omega_j(t)\le h$. For $h=1$ this means that head $j$ arrived on the cell at
step $t-1$.

\begin{lemma}[The closure is cut by $C$]\label{lem:closure}
\notformalized
Every edge of the cell network from a tainted step to an untainted step lies in $C$. Let $N_{\mathrm{run}}$ be the number of runs and
$N_{\mathrm{up}}$ the number of tainted steps of $\Out(A)$. Then:
\begin{enumerate}[label=(\alph*)]
\item every path in the cell network from a step that reads a position of $A$ to an untainted step uses an edge of $C$;
\item the set $\Gamma$ of the tails of the edges of $C$ and the tainted steps of $\Out(A)$ separates $\Out(A)$ from $\Lin(\{1,2\})$
in the dependency DAG;
\item $C$ has at most $\abs{C_0}+3N_{\mathrm{run}}$ cell edges and at most $N_{\mathrm{run}}$ state edges, and
$\abs\Gamma\le\abs{C_0}+N_{\mathrm{run}}+N_{\mathrm{up}}$.
\end{enumerate}
\end{lemma}

\begin{informalproof}
Let $p\to v$ be an edge with $p$ tainted and $v$ untainted. A state edge has $v=p+1$, so $p$ ends a run and the edge is in $C$. A cell
edge with $v\ge p+2$ is in $C_0$, since otherwise $v$ would be tainted. A wait of tape $j$ that is not in $C_0$ has $\omega_j(v)>h$,
since otherwise $v$ would be tainted; so $p$ ends a run, and the wait is in $C$. This proves the first claim, and (a): a path from a
tainted step to an untainted one has a first untainted vertex, and the edge into it leaves a tainted vertex.

For (b), take a dependency path from $\Out(v)$, $v\in A$, to a step that reads a position of $A$. If $\Out(v)$ is tainted it lies in
$\Gamma$. Otherwise the reversed path is a path in the cell network (\cref{def:cellnet}) from a tainted step to an untainted one, and
the edge into its first untainted vertex lies in $C$ and leaves a vertex of the path, which is a tail of $C$. For (c), $C_0$ has
$\abs{C_0}$ edges and tails, and every other edge of $C$ leaves the last step of a run: one state edge, none after the last step of
the trace, and at most one wait per work tape; both schedules have three.
\end{informalproof}

So it suffices to choose $C_0$ so that $\abs{C_0}$, $N_{\mathrm{run}}$ and $N_{\mathrm{up}}$ are small.

\subsection{Cascade 3}\label{app:cascade3}

Cascade~3 uses the closure of order $h=1$. Its regions are as follows. On tape~$0$, $R_r$ starts at $a_r$, with $a_1=L$ and
$a_{r+1}=a_r+2m_r$; for the backward stage, $\mathrm{IN}_r$ has size $m_r$ and $\mathrm{OUT}_r$ has size $2m_r$, with
$\mathrm{IN}_{r-1}=\mathrm{OUT}_r$. On tape~$1$, the level-$r$ scratch starts at some $s_r$, and the backward stage uses fresh regions
$\mathrm{CB}_r$ of size $m_r$. On tape~$2$, $Z_r=[z_r,z_r+m_r)$ with $z_r=L-2m_r$. Put $\Delta=\lfloor\lfloor L/3\rfloor/2\rfloor-1$. The
phases that matter are the following.

\begin{center}\small
\begin{longtable}{>{\raggedright\arraybackslash}p{0.2\textwidth}>{\raggedright\arraybackslash}p{0.74\textwidth}}
\toprule
phase & head motion \\
\midrule
\endhead
load, step $u<L$ & input head on $u$; tapes $0$ and $1$ on cell $u$; tape~$2$ on cell $L+\lceil u/2\rceil$ while $u\le\lfloor L/3\rfloor$,
then at speed~$1$ \\
rev & tape~$0$ down $[L,2L)$ and tape~$2$ down to cell $L$, tape-$2$ cell $x$ aligned with tape-$0$ cell $x+\Delta$; then tape~$2$
parks on cell~$L$; tape~$1$ idle \\
fwd $r$.transfer, step $s$ & tape~$0$ on $R_r$ offset $\min(s,m_r-1)$; tape~$1$ on scratch offset $\lceil s/2\rceil$ for
$s<2\lfloor m_r/3\rfloor$, then at speed~$1$; tape~$2$ idle \\
fwd $r$.move2 & tape~$0$ up to $R_r$ offset $2m_r-\mathrm{sh}_r$, then idle; tape~$1$ down the scratch from offset $m_r-1$ to $0$ \\
fwd $r$.copyS & tapes $0,1,2$ on $R_r$ offset $2m_r-\mathrm{sh}_r+u$, scratch offset $u$, $Z_r$ offset $u$, for $0\le u<\mathrm{sh}_r$ \\
fwd $r$.move3, copyR, move4 & tape~$0$ down to $R_r$ offset $m_r$; then $R_r$ offset $m_r+u$, scratch offset $\mathrm{sh}_r+u$, $Z_r$
offset $\mathrm{sh}_r+u$; then tape~$0$ up to $R_{r+1}$ offset $m_r-1$ \\
fwd $r$.back, move5 & tape~$0$ down $R_{r+1}$ and tape~$2$ down $Z_r$, offset aligned with offset; then tape~$2$ up through $Z_r$ to
$z_{r+1}$ \\
comb $r$.move & tape~$0$ up $\mathrm{IN}_r$ from offset~$0$; tape~$1$ up from $\mathrm{CB}_{r+1}$ offset $\mathrm{sh}_{r+1}$ (from the
end of the scratch if $r=r_{\max}$) to $\mathrm{CB}_r$ offset $m_r-1$; $\mathrm{IN}_r$ offset $s$ is aligned with $\mathrm{CB}_{r+1}$
offset $\mathrm{sh}_{r+1}+s$ while $s<m_{r+1}-\mathrm{sh}_{r+1}$, and with $\mathrm{CB}_r$ offset $s-c_r$, where
$c_r=m_{r+1}+2-\mathrm{sh}_{r+1}$, once $s\ge c_r$; then tape~$2$ moves to $z_r+m_r-1$: one step from $z_{r+1}$ if $r<r_{\max}$, and up
through the offsets $0,\dots,m_r-2$ of $Z_{r_{\max}}$ if $r=r_{\max}$ \\
comb $r$.blocks & $2^r$ blocks of $m_r/2^r$ offsets, from the top: all three heads down together over $\mathrm{IN}_r$, $\mathrm{CB}_r$
and $Z_r$ at equal offsets; then tape~$1$ up and down across the block while tapes $0$ and $2$ wait \\
comb $r$.move2 & tape~$0$ up through $\mathrm{IN}_r$ to $\mathrm{OUT}_r$; tapes $1$ and $2$ idle \\
comb $r$.dbl & $\mathrm{CB}_r$ offset $z$ aligned with $\mathrm{OUT}_r$ offsets $2z$ and $2z+1$ \\
comb $r$.move3, shift, move4, rem & tape~$1$ alone down to offset $\mathrm{sh}_r-1$; then $\mathrm{OUT}_r$ offset $y$ aligned with
$\mathrm{CB}_r$ offset $(\lfloor y/2\rfloor+\mathrm{sh}_r)\bmod m_r$, both moving down, with tape~$1$ alone moving up between the two parts \\
rewind & tape~$1$ down from $\mathrm{CB}_1$ offset $\mathrm{sh}_1$ to cell~$0$; tapes $0$ and $2$ idle \\
emit, step $u$ & tape~$0$ on $\mathrm{OUT}_1$ offset $u$; tape~$1$ on cell $u$; tape~$2$ on cell $\min(u,\sum_rm_r-1)$; output bit $u$
appended \\
\bottomrule
\end{longtable}
\end{center}

The positioning phases omitted from the table (the initial positioning, the moves into and out of rev, and the move that starts each
forward level) read no live tainted cell. We use three properties of the trace.
\begin{enumerate}[label=(\roman*)]
\item Every phase is a \emph{segment}, or a few of them: an interval of steps in which every head is either idle throughout, or
monotone and waiting at most one step at a time. To get segments, cut rev, transfer, move2, comb move and emit where a head stops early,
cut each block into its sweeps, and make the last step of copyS, at which all three heads wait, a segment of its own.
\item After the current level, the regions of tape~$0$ below the current one, the scratch regions, and $\mathrm{CB}_{r'}$ for
$r'>r+1$ are not visited again, except by tape~$1$ in the rewind. $Z_r$ is visited after move5 only in comb $r$.blocks and in the emit.
The exceptions are single idle cells: offset~$0$ of $Z_r$ while comb $r$.move2 to rem run and again in comb $(r-1)$.move, and offset~$0$
of $Z_1$ and of $\mathrm{OUT}_1$ during the rewind. Each adds at most one run of at most two steps.
\item In every segment at most one head moves at half speed (in load, transfer, dbl, shift and rem).
\end{enumerate}

\begin{lemma}[Segments]\label{lem:c3-segments}
\notformalized
In a segment of Cascade~3, every tainted step is (1) a step at which a head that is not idle scans a live tainted cell, (2) one of the
first two steps of an idle head on a live tainted cell, (3) a step that reads a position of $A$, or (4) the step right after one of
these, or the step after that when the first is a short wait of an idle head. Hence the number of runs in a segment is at most the
number of maximal intervals of live tainted cells that it reads, over all tapes, plus the number of runs of reads of $A$; and each run
writes one interval of cells on each tape whose head moves.
\end{lemma}

\begin{informalproof}
A monotone head stays on an interval of cells during an interval of steps. Its first step on a cell is a revisit, or at the start of
the segment a wait of either kind, and its second step there, if any, is a short wait. A short wait at the first step $t_0$ of the
segment from a tainted step $t_0-1$ reads a cell that is live tainted at $t_0$, so $t_0$ is of kind~(1). An idle head's first two steps
are a revisit and a short wait, and the rest are waits that are not short. Cells written inside the segment are revisited inside it only
by waits. So a tainted step of none of the kinds (1)--(3) is tainted through a short wait from the tainted step before it. By~(iii), all
short waits of moving heads in the segment belong to one head $j$. After its wait at step $t$, head $j$ moves, so step $t+1$ is no wait
of $j$, and no wait of any other moving head either; so the extension stops after one step. An idle head's short wait at $t_0+1$ can hand
over to a short wait of $j$ at $t_0+2$, so the extension stops after at most two steps. It is adjacent to its run, so it adds no run and
extends each written interval by at most two cells.
\end{informalproof}

The cut is placed at a level $\rho$ with $1\le\rho\le r_{\max}$: $C_0$ is the set of tape-$1$ edges that leave tainted steps of
fwd $\rho$.transfer and whose next visit lies after the transfer. The waits of tape~$1$ inside the transfer are not cut. This is well
defined: no edge of $C_0$ ends at or before the end of the transfer, so the taint up to that point does not depend on $C_0$.

\begin{lemma}[Forward windows]\label{lem:c3-forward}
\notformalized
For $1\le r\le\rho$ let $W_r$ be the set of offsets of $R_r$ that are live tainted when fwd $r$.transfer starts. Then
$W_1=\{\lceil u/2\rceil+\Delta:u\in A\}$ and $W_{r+1}=\{\lceil s/2\rceil:s\in W_r\}$ for $r<\rho$. Each $W_r$ is an interval with
\[
  \abs{W_r}\le\frac{db}{2^{r-1}}+2,\qquad \frac{m_r}3-2<\min W_r\le\max W_r\le\frac{m_r}3+\frac{3db}{2^r}+1 .
\]
Moreover $\abs{C_0}\le db/2^\rho+2$, and the forward stage up to the end of fwd $\rho$.transfer has $O(\rho)$ runs.
\end{lemma}

\begin{informalproof}
The bounds hold for $W_1$: the interval $\{\lceil u/2\rceil:u\in A\}$ has at most $db+1$ elements, from $\lceil db/2\rceil$ to
$\lceil(3db-1)/2\rceil$, and $m_1/3-2\le\Delta\le m_1/3-1$. They pass from $W_r$ to $W_{r+1}$, because $\lceil s/2\rceil\in[s/2,(s+1)/2]$
and the image of an interval of $n$ integers has at most $n/2+1$ elements. Since $3db/2^r=3m_r/d\le3m_r/32$ and $m_r\ge64$,
\begin{equation}\label{eq:c3-window}
  \max W_r\le\tfrac{41}{96}m_r+1<2\lfloor m_r/3\rfloor,\qquad \max\{\lceil s/2\rceil:s\in W_r\}\le\tfrac{41}{192}m_r+1<\mathrm{sh}_r-2 .
\end{equation}

Load: the tainted steps are the reads of $A$ and, at most, the step after the last one, a short wait of tape~$2$. They write tape-$0$
cells, which are not visited again; the loaded copy of rows~$1$ and~$2$ on tape~$1$, cells $[db,3db]$, which is live until the rewind;
and the tape-$2$ cells $L+\lceil u/2\rceil$, $u\in A$, as $3db<\lfloor L/3\rfloor$. Rev reads these tape-$2$ cells and writes the tape-$0$
cells $L+\lceil u/2\rceil+\Delta$, that is, $W_1$; tape~$1$ is idle, so its wait is cut, and the tape-$2$ cells above $L$ are not
visited again.

In fwd $r$.transfer the steps $s\in W_r$ are tainted, and so is $s+1$ for odd $s$, by a short wait of tape~$1$. By \eqref{eq:c3-window}
tape~$1$ is in its half-speed part, and the scratch cells it writes are the offsets $S_r=\{\lceil s/2\rceil:s\in W_r\}$. For $r=\rho$ the
edges leaving them are $C_0$, so $\abs{C_0}\le\abs{S_\rho}\le\abs{W_\rho}/2+1$. For $r<\rho$, the scratch window $S_r$ lies in
$[0,\mathrm{sh}_r-2)$. It is read in move2, where tape~$0$ is already idle because $m_r-1-\max S_r>m_r-\mathrm{sh}_r+1$; in copyS, aligned
with $Z_r$ offset $u$ and with $R_r$ offset $2m_r-\mathrm{sh}_r+u$; and next in the rewind. The copy on $R_r$ is read in move3 and move4
and then never again. The copy on $Z_r$ is read in back, which writes the offsets $S_r$ of $R_{r+1}$, so $W_{r+1}=S_r$, and in move5.
CopyR touches only scratch and $Z_r$ offsets at least $\mathrm{sh}_r$ and $R_r$ offsets below $2m_r-\mathrm{sh}_r$, none of them tainted.
Each of transfer, move2, copyS, move3, move4, back and move5 reads at most two intervals of live tainted cells, so it has at most two
runs (\cref{lem:c3-segments}).
\end{informalproof}

\begin{lemma}[Clean middle]\label{lem:c3-middle}
\notformalized
No step of fwd $r$ for $r>\rho$, of fwd $\rho$ after its transfer, or of comb $r$ for $r\ge\rho$ is tainted.
\end{lemma}

\begin{informalproof}
The cut $C_0$ blocks the scratch window of level $\rho$, so the steps of move2 that read it are untainted and rewrite it with
untainted contents. $Z_\rho$ and $R_{\rho+1}$ are therefore written by untainted steps, and so is everything built from them. These
phases visit only regions of levels at least $\rho$, by~(ii), and the tainted windows of lower levels ($Z_{r'}$, the scratch regions,
$\mathrm{CB}_{r'}$ and the loaded copy) lie in regions they do not visit. The one tainted window of level $\rho$ is
$W_\rho\subseteq R_\rho$; after the transfer, tape~$0$ visits only offsets of $R_\rho$ at least $m_\rho-1$, and $W_\rho\subseteq[0,m_\rho-1)$.
\end{informalproof}

For $r<\rho$ let $U_r$ be the set of offsets $y\in[0,m_r)$ at which $\mathrm{IN}_r$, $Z_r$ or $\mathrm{CB}_r$ is live tainted when
comb $r$.blocks starts, $u_r$ its number of maximal intervals and $\mu_r=\abs{U_r}$; let $h_r=m_r/2^r$.

\begin{lemma}[Backward stage]\label{lem:c3-backward}
\notformalized
We have $u_{\rho-1}\le1$ and $\mu_{\rho-1}\le db/2^{\rho-1}+2$, and for $r<\rho-1$
\[
  u_r\le6u_{r+1}+8,\qquad \mu_r\le10\mu_{r+1}+\frac{db}{2^r}+O(u_{r+1}).
\]
Comb $r$ has at most $c\,(u_r+u_{r+1}+\mu_r/h_r+1)$ runs, for an absolute constant $c$.
\end{lemma}

\begin{informalproof}
Comb $\rho-1$ reads $\mathrm{IN}_{\rho-1}$ and the comb region of level $\rho$, which come from comb $\rho$ and are clean
(\cref{lem:c3-middle}); so $U_{\rho-1}$ is the window $S_{\rho-1}$ of $Z_{\rho-1}$, one interval with at most $db/2^{\rho-1}+2$ cells.

In comb $r$, blocks moves the three heads in alignment, so afterwards all three regions are tainted exactly on $U_r$, together with
one cell at the bottom of a block, where the waiting heads pass taint into the first two steps of the excursion; the excursions only
reread tape-$1$ cells of the same block. Move2 rereads $\mathrm{IN}_r$, which is not visited afterwards. Dbl taints $\mathrm{OUT}_r$ on
$2U_r$; shift and rem, through $y\mapsto(\lfloor y/2\rfloor+\mathrm{sh}_r)\bmod m_r$, add $2(U_r-\mathrm{sh}_r)$ to $\mathrm{OUT}_r$
and $U_r+\mathrm{sh}_r$ to $\mathrm{CB}_r$. So $\mathrm{IN}_{r-1}=\mathrm{OUT}_r$ and $\mathrm{CB}_r$ each have at most $2u_r+2$
intervals, of total size at most $4\mu_r$ and $2\mu_r$.

In comb $(r-1)$.move the alignments of the table add to $\mathrm{IN}_{r-1}$ the image of the comb region of level $r$ (its offsets at
least $\mathrm{sh}_r$, shifted by $\mathrm{sh}_r$), and to $\mathrm{CB}_{r-1}$ the image of the offsets of $\mathrm{IN}_{r-1}$ at least
$c_{r-1}$. The image of $\mathrm{CB}_r$ lands at offsets of $\mathrm{IN}_{r-1}$ below $m_r-\mathrm{sh}_r<c_{r-1}$, where tape~$1$ is still
in the old comb region, so the steps that read it write no cell of $\mathrm{CB}_{r-1}$, and $\mathrm{CB}_{r-1}$ receives at most the
$2u_r+2$ intervals of $\mathrm{OUT}_r$. Hence $u_{r-1}\le(2u_r+2)+(2u_r+2)+2+(2u_r+2)$ and
$\mu_{r-1}\le4\mu_r+2\mu_r+\abs{S_{r-1}}+4\mu_r$, up to $O(1)$ per interval.

Runs: each segment of comb $r$ has at most as many runs as intervals it reads (\cref{lem:c3-segments}). In blocks an interval is read
in every block it meets, and the $u_r$ intervals meet at most $\mu_r/h_r+2u_r$ blocks, three sweeps each. Comb $r$.move reads
$\mathrm{OUT}_{r+1}$ and $\mathrm{CB}_{r+1}$, which gives the term $u_{r+1}$.
\end{informalproof}

Solving the recurrences gives $u_r\le4\cdot6^{\rho-1-r}$. For the sizes, $\kappa_r=\mu_r/(db)$ satisfies $\kappa_{r-1}\le10\kappa_r+2^{1-r}$
up to the terms per interval, so
\[
  \kappa_r\le\sum_{j=r}^{\rho-1}10^{j-r}2^{1-j}=2^{1-r}\sum_{j=r}^{\rho-1}5^{j-r}\le\tfrac54\,2^{1-r}\,5^{\rho-r},
\]
a bound geometric in $r$; the terms per interval add $O(\sum_j10^{j-r}6^{\rho-j})=O(10^{\rho-r})$ cells at level $r$. Hence
$\sum_{r<\rho}\mu_r\le c\,(5^\rho db+10^\rho)$, $\sum_{r<\rho}\mu_r/h_r=\sum_r\mu_r4^r/L\le c\,(5^\rho/d+40^\rho/L)$, and the backward
stage has at most $c\,(6^\rho+5^\rho/d+40^\rho/L)$ runs.

\paragraph{Rewind and emit.} The rewind moves only tape~$1$, down from $\mathrm{CB}_1$ offset $\mathrm{sh}_1$. It reads the live tainted
intervals of tape~$1$: the comb regions and the scratch windows of levels below $\rho$, and the loaded copy $[db,3db]$; so it has at most
$c\,6^\rho+2\rho+2$ runs. Afterwards only the loaded copy among them is visited again. The emit's step $u$ reads $\mathrm{OUT}_1$
offset $u$, tape-$1$ cell $u$ and tape-$2$ cell $\min(u,\sum_rm_r-1)$, and appends output bit $u$. A tainted step of the emit reads a
tainted cell of the loaded copy, at $u\in[db,3db]$; of $\mathrm{OUT}_1$, which has at most $2u_1+2$ intervals and size at most $4\mu_1$;
or of $Z_r$ for $r<\rho$, which is tainted on $U_r$. Among the outputs $u\in[db,3db]$, those of $\Out(A)$ are the $4b$ bits of the
entries $(1,1)$, $(1,2)$, $(2,1)$ and $(2,2)$. Elsewhere, the outputs of $\Out(A)$ form one run of $2b$ positions in each block of $db$
output positions, so an interval of length $\lambda$ contains at most $2b(\lambda/(db)+2)$ of them. The other tainted intervals of the
emit number at most $c\,6^\rho$ and have total length at most $4\mu_1+\sum_r\mu_r\le c\,(5^\rho db+10^\rho)$. Hence
\[
  N_{\mathrm{up}}\le4b+2b\bigl(c\,5^\rho+c\,10^\rho/(db)+2c\,6^\rho\bigr)\le4b+c'\,6^\rho b\qquad\text{once }(5/3)^\rho\le db .
\]
Altogether, for every $1\le\rho\le r_{\max}$,
\begin{equation}\label{eq:c3-total}
  \abs{C_0}\le\frac{db}{2^\rho}+2,\qquad N_{\mathrm{run}}\le c''\Bigl(6^\rho+\frac{5^\rho}d+\frac{40^\rho}L+\rho\Bigr),\qquad
  N_{\mathrm{up}}\le4b+c'\,6^\rho b .
\end{equation}

\subsection{Cascade 4}\label{app:cascade4}

Cascade~4 uses the closure of order $h=2$, because its final sweep moves tapes~$0$ and~$1$ once every three steps. All its regions are
disjoint, separated by two buffer cells. On tape~$2$: the frames $Z_r=[z_r,z_r+m_r)$ and $\mathrm{BS}_r=[z_r+m_r+2,z_r+3m_r+2)$ for
$r=1,\dots,r_{\max}$, with $z_1=0$ and $z_{r+1}=z_r+3m_r+4$, ending at $E_2=\sum_{r\le r_{\max}}(3m_r+4)<3L$; then the seed region
$\mathrm{SD}=[E_2+2,E_2+2+L)$. On tape~$0$: the loaded copy $\mathrm{LD}_0=[0,L)$; $X_1=[L+2,2L+2)$; for $r=1,\dots,r_{\max}$, a region
$\mathrm{H2R}_r$ of size $m_r$ and the region $X_{r+1}$ of size $m_r=2m_{r+1}$, the input of level $r+1$; then, for $r=r_{\max}$ down to
$1$, regions $\mathrm{CM}_r$ of size $m_r$, $\mathrm{U2R}_r$ of size $2m_r$ and $B_{r-1}$ of size $2m_r$, with $B_{r_{\max}}=X_{r_{\max}+1}$
and $B_0$ of size $L$. On tape~$1$: the loaded copy $\mathrm{LD}_1=[0,L)$; for $r=1,\dots,r_{\max}$, regions $\mathrm{H1}_r$ and
$\mathrm{H2}_r$ of size $m_r$; then, for $r=r_{\max}$ down to $1$, regions $\mathrm{MZ}_r$ of size $m_r$ and $\mathrm{U1}_r$, $\mathrm{U2}_r$
of size $2m_r$. In the table, $Y[x]$ is offset $x$ of region $Y$.

\begin{center}\small
\begin{longtable}{>{\raggedright\arraybackslash}p{0.2\textwidth}>{\raggedright\arraybackslash}p{0.74\textwidth}}
\toprule
phase & head motion \\
\midrule
\endhead
initial positioning, load, rev & tape~$2$ up from $0$ to $\mathrm{SD}[0]$ while the input head stays on bit~$0$; load, step $u<L$:
input head on $u$, tapes $0$ and $1$ on $\mathrm{LD}_0[u]$ and $\mathrm{LD}_1[u]$, tape~$2$ on $\mathrm{SD}[\lceil u/2\rceil]$ while
$u\le\lfloor L/3\rfloor$, then at speed~$1$; tape~$0$ up to $X_1[L-1]$; rev: $X_1[x+\Delta]$ aligned with $\mathrm{SD}[x]$, then tape~$2$
parks on $\mathrm{SD}[0]$; tape~$2$ down to $z_1=0$ \\
fwd $r$.move, transfer & tape~$0$ idle on $X_r[0]$, tape~$1$ up to $\mathrm{H1}_r[0]$; transfer, step $s$: $X_r[\min(s,m_r-1)]$ aligned
with $\mathrm{H1}_r[\lceil s/2\rceil]$ for $s<2\lfloor m_r/3\rfloor$, then at speed~$1$, ending with one step at which no head moves \\
fwd $r$.shift, rev2 & tape~$0$ reads $X_r[m_r],\dots,X_r[2m_r-1]$, cyclically shifted by $\mathrm{sh}_r$, into the fresh $\mathrm{H2}_r$;
then $\mathrm{H2}_r$ is reversed into the fresh $\mathrm{H2R}_r$; tape~$2$ idle on $z_r$ \\
fwd $r$.merge.move, merge, back, move5 & merge.move, $m_r+2$ steps: tape~$0$ down $\mathrm{H2R}_r$ from offset $m_r-1$ to $0$ ($m_r-1$
steps, then it parks), tape~$1$ down from $\mathrm{H2}_r[0]$ to $\mathrm{H1}_r[0]$, so that $\mathrm{H1}_r[x]$ is aligned with
$\mathrm{H2R}_r[x-3]$, tape~$2$ idle on $z_r$; merge: $\mathrm{H2R}_r[u]$ and $\mathrm{H1}_r[u]$ into $Z_r[u]$; back: $Z_r[u]$ into
$X_{r+1}[u]$, downwards; move5 ($r<r_{\max}$): tape~$2$ up from $z_r$ to $z_{r+1}$ across $Z_r$ and the clean $\mathrm{BS}_r$ \\
comb $r$.1a.move, 1a & 1a.move: tape~$0$ up $B_r[0..m_r-1]$; tape~$1$ up across $\mathrm{U2}_{r+1}$ (or $\mathrm{H2}_{r_{\max}}$) and
the fresh $\mathrm{MZ}_r$; tape~$2$ down from $\mathrm{BS}_{r+1}[0]$ across $Z_{r+1}$ and $\mathrm{BS}_r$ to $Z_r[m_r-1]$ (up from
$Z_{r_{\max}}[0]$ if $r=r_{\max}$); 1a: $B_r[u]$ and $Z_r[u]$ into $\mathrm{MZ}_r[u]$, downwards \\
comb $r$.blocks.move, blocks & blocks.move: tape~$0$ up across $B_r$ into the fresh $\mathrm{CM}_r[m_r-1]$, tape~$1$ up $\mathrm{MZ}_r$;
blocks: $2^r$ blocks of height $h_r=m_r/2^r$, from the top; in each, tape~$1$ down $h_r-1$ cells, up $h_r-1$ cells, then tapes $0$
and $1$ down $h_r$ cells together \\
comb $r$.dbl, shift, rev2 & dbl: $\mathrm{CM}_r[z]$ into $\mathrm{U1}_r[2z]$ and $\mathrm{U1}_r[2z+1]$; shift:
$\mathrm{CM}_r[(z+\mathrm{sh}_r)\bmod m_r]$ into $\mathrm{U2}_r[2z]$ and $\mathrm{U2}_r[2z+1]$; rev2: $\mathrm{U2}_r$ reversed into
$\mathrm{U2R}_r$; tape~$2$ idle on $z_r$ \\
comb $r$.merge.move, merge, out & merge.move, $2m_r+2$ steps: tape~$0$ down $\mathrm{U2R}_r$, tape~$1$ down from $\mathrm{U2}_r[0]$ to
$\mathrm{U1}_r[0]$, tape~$2$ up from $Z_r[0]$ to $\mathrm{BS}_r[0]$ ($m_r+2$ steps); merge: $\mathrm{U2R}_r[u]$ and $\mathrm{U1}_r[u]$ into
$\mathrm{BS}_r[u]$; out: $\mathrm{BS}_r[u]$ into $B_{r-1}[u]$, downwards \\
rewind, emit & rewind: tape~$0$ idle on $B_0[0]$, tape~$1$ down from $\mathrm{U1}_1[L-1]$ to cell~$0$, tape~$2$ down from
$\mathrm{BS}_1[0]$ to cell~$0$; emit, steps $3u+i$ for $u<L$ and $i=0,1,2$: tape~$0$ on $B_0[u]$ and tape~$1$ on $\mathrm{LD}_1[u]$, both
moving at $i=2$, tape~$2$ on cell $\min(3u+i,E_2-1)$, output bit $u$ appended at $i=2$ \\
\bottomrule
\end{longtable}
\end{center}

Split each phase into segments at the steps where a head stops early or where no head moves, and each block into its sweeps; in
particular split fwd $r$.transfer where tape~$0$ stops while tape~$1$ still moves, and the emit where tape~$2$ parks on cell
$E_2-1$. Then in every segment each head is idle or monotone, waiting at most two consecutive steps on a cell, and whenever more than one
moving head waits, which happens only in the emit, the moving heads that wait all wait at the same steps.

\begin{lemma}[Segments of Cascade~4]\label{lem:c4-segments}
\notformalized
In a segment $[t_0,t_1)$ of Cascade~4, the number of runs is at most the number of maximal intervals of live tainted cells read by the
moving heads, plus the number of idle heads on a live tainted cell that arrived there at most two steps before $t_0$, plus the number
of runs of reads of $A$; and each run writes at most one interval of cells on each tape whose head moves.
\end{lemma}

\begin{informalproof}
As in \cref{lem:c3-segments}, a tainted step that neither reads a live tainted cell with a moving head nor is one of the first steps
of an idle head on a live tainted cell nor reads $A$ is tainted through a short wait from the step before it. An idle head that parked
three or more steps before $t_0$ has wait age at least $3$ throughout the segment and passes no taint. The moving heads that wait do so
at the same steps, at most two in a row, and then all advance together with the heads that do not wait; so a chain of short waits ends
within two steps. The one longer chain, of at most five tainted steps, occurs when an idle head's short waits at $t_0+1,t_0+2$ hand over
to a moving head's waits at $t_0+3,t_0+4$, at the step where tape~$2$ parks in the emit; it stays adjacent to its run. So each extension
adds no run and only boundedly many written cells.
\end{informalproof}

The cut is placed at a level $\rho$ with $1\le\rho\le r_{\max}$: $C_0$ is the set of tape-$1$ edges that leave tainted steps of
fwd $\rho$.transfer and whose next visit lies after the transfer; as for Cascade~3, this is well defined.

\begin{lemma}[Forward windows of Cascade~4]\label{lem:c4-forward}
\notformalized
For $1\le r\le\rho$ let $W_r$ be the set of offsets of $X_r$ that are live tainted when fwd $r$.transfer starts, and put
$S_r=\{\lceil s/2\rceil:s\in W_r\}$. Then $W_1=\{\lceil u/2\rceil+\Delta:u\in A\}$ and $W_{r+1}=S_r\cup(S_r-3)$ for $r<\rho$. Each $W_r$
is an interval of at least $6$ offsets with
\[
  \abs{W_r}\le\frac{db}{2^{r-1}}+8,\qquad \frac{m_r}3-6\le\min W_r\le\max W_r\le\frac{m_r}3+\frac{3db}{2^r}+1\le2\lfloor m_r/3\rfloor-3 ,
\]
and $S_r\subseteq[7,\mathrm{sh}_r-3]$. Moreover $\abs{C_0}=\abs{S_\rho}\le db/2^\rho+5$, and the forward stage up to the end of
fwd $\rho$.transfer has $O(\rho)$ runs.
\end{lemma}

\begin{informalproof}
Load writes the tainted cells $\mathrm{SD}[\lceil u/2\rceil]$, $u\in A$, since $3db\le\lfloor L/3\rfloor$, together with the loaded
copy $\mathrm{LD}_1[db..3db]$; rev aligns $\mathrm{SD}[x]$ with $X_1[x+\Delta]$, with no waits of the moving heads and tape~$1$ parked for
more than two steps, which gives $W_1$. The bounds pass from $W_r$ to $W_{r+1}$ as in \cref{lem:c3-forward}. If $K$ is an interval of
at least $3$ offsets, $(K-3)\cup K$ is the interval $[\min K-3,\max K]$; since $\abs{S_r}\ge\abs{W_r}/2\ge3$, each $W_{r+1}$ is again one
interval, of at least $6$ offsets, with $\abs{W_{r+1}}\le\abs{W_r}/2+4$, $\min W_{r+1}=\min S_r-3\ge m_{r+1}/3-6$ and
$\max W_{r+1}=\max S_r\le m_{r+1}/3+3db/2^{r+1}+1$. As in \eqref{eq:c3-window}, $\max W_r\le\frac{41}{96}m_r+1\le2\lfloor m_r/3\rfloor-3$,
and $S_r\subseteq[\frac{m_r}6-3,\frac{41}{192}m_r+1]\subseteq[7,\mathrm{sh}_r-3]$, since $m_r\ge64$.

So fwd $r$.transfer stays in its half-speed part on $W_r$ and writes the tainted offsets $S_r$ of $\mathrm{H1}_r$; for $r=\rho$ the
edges leaving them are $C_0$. As $W_r\subseteq[0,m_r-2]$, the second half of $X_r$ is untainted, so shift and rev2 have no tainted step,
and $\mathrm{H2}_r$ and $\mathrm{H2R}_r$ are clean. In fwd $r$.merge.move, tape~$1$ visits $\mathrm{H1}_r[x]$, $x\in S_r$, at the same
step at which tape~$0$ visits $\mathrm{H2R}_r[x-3]$, with no waits of either: tape~$0$ parks on $\mathrm{H2R}_r[0]$ only at
step $m_r-1$ of merge.move, when tape~$1$ is at offset $3$ of $\mathrm{H1}_r$, below $\min S_r\ge7$. Tape~$2$ has been parked on $z_r$ since
before the transfer. So merge taints $Z_r$ exactly at the offsets $S_r\cup(S_r-3)$, and back copies $Z_r$ onto $X_{r+1}$, which gives $W_{r+1}$;
move5 crosses this window and leaves $\mathrm{BS}_r$ untainted. Each forward level has a bounded number of segments, each reading at
most two intervals, which gives the $O(\rho)$ runs.
\end{informalproof}

\begin{lemma}[Clean middle of Cascade~4]\label{lem:c4-middle}
\notformalized
No step of fwd $\rho$ after its transfer, of fwd $r$ for $r>\rho$, or of comb $r$ for $r\ge\rho$ is tainted.
\end{lemma}

\begin{informalproof}
By \cref{lem:c4-forward} the second half of $X_\rho$ and $\mathrm{H1}_\rho[m_\rho-1]$ are untainted, and $C_0$ cuts every tape-$1$ edge
from a tainted step of fwd $\rho$.transfer to the later phases. From the end of that transfer to the end of comb $\rho$.out, tape~$0$
visits only cells at or beyond $X_\rho[m_\rho-1]$, tape~$1$ only cells at or beyond $\mathrm{H1}_\rho$, and tape~$2$ only cells of
$[z_\rho,E_2)$. The live tainted regions of levels at most $\rho$ (on tape~$0$, $\mathrm{LD}_0$, the first halves of $X_{r'}$ for
$r'\le\rho$ and $\mathrm{H2R}_{r'}$ for $r'<\rho$; on tape~$1$, $\mathrm{LD}_1$, $\mathrm{H1}_{r'}$ for $r'<\rho$ and the cut cells of
$\mathrm{H1}_\rho$; on tape~$2$, $\mathrm{SD}$ and $Z_{r'}$ for $r'<\rho$) are not reached by an uncut edge in these phases.
\end{informalproof}

For $1\le r\le\rho-1$ let $U_r\subseteq[0,m_r)$ be the set of offsets at which $B_r$ or $Z_r$ is live tainted when comb $r$.1a starts,
with $u_r$ maximal intervals and $\mu_r=\abs{U_r}$, and let $V_r\subseteq[0,2m_r)$ be the set of live tainted offsets of $\mathrm{BS}_r$
at that moment, with $v_r$ maximal intervals and $\eta_r=\abs{V_r}$. For $r\le\rho-1\le k-2$ the block height is
$h_r=m_r/2^r=d^2b/4^r\ge16b\ge16$.

\begin{lemma}[Backward stage of Cascade~4]\label{lem:c4-backward}
\notformalized
Let $\rho\ge2$. For $1\le r\le\rho-1$, $V_r\subseteq[\frac32m_r+2,2m_r)$; $u_{\rho-1}\le1$, $\mu_{\rho-1}\le db/2^{\rho-1}+8$ and
$V_{\rho-1}=\varnothing$; and for $1<r\le\rho-1$
\begin{gather*}
  u_{r-1}\le10u_r+v_r+19,\qquad v_{r-1}\le7u_r+v_r+13,\\
  \mu_{r-1}\le17\mu_r+\eta_r+\frac{db}{2^{r-1}}+28u_r+30,\qquad \eta_{r-1}\le12\mu_r+\eta_r+20u_r+12 .
\end{gather*}
Comb $r$ has $O(u_r+v_r+u_{r+1}+v_{r+1}+\mu_r/h_r+1)$ runs, with $u_\rho=v_\rho=0$, and leaves $Z_r$ and $\mathrm{BS}_r$ with
$O(u_r+v_r+1)$ tainted intervals and $O(\mu_r+\eta_r+u_r+1)$ tainted cells.
\end{lemma}

\begin{informalproof}
At $r=\rho-1$, comb $\rho$ is clean (\cref{lem:c4-middle}), so comb $(\rho-1)$.1a.move crosses only clean cells and has no tainted step;
$U_{\rho-1}$ is the window $W_\rho$ that merge wrote on $Z_{\rho-1}$, and $V_{\rho-1}=\varnothing$. Now follow comb $r$ for $r\le\rho-1$,
assuming $V_r\subseteq[\frac32m_r+2,2m_r)$.
\begin{enumerate}[label=(\arabic*)]
\item 1a writes a subset of $U_r$ onto the fresh $\mathrm{MZ}_r$. In blocks.move tape~$0$ crosses $B_r$ and tape~$1$ crosses
$\mathrm{MZ}_r$; a short wait of tape~$2$ on $Z_r[0]$ adds at most offset $1$. Tape~$0$ then crosses two clean cells into the fresh
$\mathrm{CM}_r$ while tapes $1$ and~$2$ have wait age at least $3$, so $\mathrm{CM}_r$ enters blocks clean.
\item In each block $[B,B+h_r)$ tape~$1$ turns on $B$ without waiting, and the third sweep ends by stepping onto the next block with
wait age $0$ on both moving heads, so no taint spills across blocks. If $B+h_r-1\in U_r$, the two-step short wait of tape~$0$ at the top
of the block taints at most $\mathrm{MZ}_r[B+h_r-2]$ and $\mathrm{MZ}_r[B+h_r-3]$; otherwise tape~$0$ reaches the third sweep with wait
age $2h_r-2>2$. So the tainted cells of $\mathrm{CM}_r$ form at most $u_r$ intervals with at most $\mu_r+2u_r$ cells, and blocks has
$O(\mu_r/h_r+u_r)$ runs.
\item Dbl writes $\mathrm{U1}_r$ with at most $u_r$ intervals and $2\mu_r+4u_r$ cells; shift and rev2 write $\mathrm{U2}_r$ and
$\mathrm{U2R}_r$ with at most $u_r+3$ intervals and $2\mu_r+4u_r+2$ cells.
\item In merge.move, tapes $0$, $1$ and $2$ cross $\mathrm{U2R}_r$, $\mathrm{U1}_r$ and $Z_r$ together, and tape~$2$ then parks on the
clean $\mathrm{BS}_r[0]$, as $0\notin V_r$. Afterwards $Z_r$ has at most $3u_r+4$ tainted intervals and $5\mu_r+8u_r+3$ tainted cells.
Merge unites $\mathrm{U2R}_r$, $\mathrm{U1}_r$ and $V_r$ on $\mathrm{BS}_r$, and out copies $\mathrm{BS}_r$ onto $B_{r-1}$: at most
$6u_r+v_r+9$ intervals and $10\mu_r+\eta_r+16u_r+8$ cells on each.
\item For $r>1$, in comb $(r-1)$.1a.move, $5m_r+7$ steps, tape~$0$ crosses $B_{r-1}[0..2m_r-1]$, tape~$1$ crosses $\mathrm{U2}_r$ and
enters the fresh $\mathrm{MZ}_{r-1}$, and tape~$2$ crosses two buffer cells, then $Z_r$, and enters the clean $\mathrm{BS}_{r-1}$ at step
$m_r+5$. The buffer cells $z_r+m_r$ and $z_r+m_r+1$ were also crossed by tape~$2$ in comb $r$.merge.move, so they may be tainted; they
meet $B_{r-1}[1..2]$ and add at most one interval and two cells, which the constants $19$ and $30$ of the lemma absorb. From step $2m_r+3$ on, tape~$0$ has wait age at least $4$, and tapes $1$ and $2$ scan only clean cells. So $\mathrm{MZ}_{r-1}$
stays clean, and $\mathrm{BS}_{r-1}$ is tainted only during the steps $[m_r+5,2m_r+2]$, which reach its offsets in
$[3m_r+2,4m_r)=[\frac32m_{r-1}+2,2m_{r-1})$. Counting the intervals and cells read in these steps gives the bounds on $v_{r-1}$ and
$\eta_{r-1}$; counting those read in the steps $[0,2m_r-1]$, together with the window $W_r$ of $Z_{r-1}$ (one interval, at most
$db/2^{r-1}+8$ cells), gives the bounds on $u_{r-1}$ and $\mu_{r-1}$.\qedhere
\end{enumerate}
\end{informalproof}

To solve the recurrences, put $\hat u_r=7u_r+v_r$ and $\hat\mu_r=12\mu_r+\eta_r$. Then $\hat u_{\rho-1}\le7$ and
$\hat u_{r-1}\le77u_r+8v_r+146\le11\hat u_r+146$, so $\hat u_r+14.6\le21.6\cdot11^{\rho-1-r}$ by induction, and
$\hat u_r\le22\cdot11^{\rho-1-r}$. Likewise $\hat\mu_{\rho-1}\le12db/2^{\rho-1}+96$ and
$\hat\mu_{r-1}\le216\mu_r+13\eta_r+12db/2^{r-1}+356u_r+372\le18\hat\mu_r+12db/2^{r-1}+51\hat u_r+372$, whence
\[
  \hat\mu_r\le c\,\bigl(2^{-r}9^{\rho-r}db+18^{\rho-r}\bigr),\qquad \sum_{r<\rho}(\mu_r+\eta_r)\le c\,(9^\rho db+18^\rho),\qquad
  \sum_{r<\rho}\frac{\mu_r}{h_r}\le c\,\Bigl(\frac{9^\rho}d+\frac{18^\rho}{d^2b}\Bigr),
\]
and the backward stage has $O(11^\rho+9^\rho/d+18^\rho/L)$ runs.

\paragraph{Rewind and emit.} In the rewind, tape~$0$ stays on $B_0[0]$ while tapes $1$ and $2$ move down to cell $0$, in two segments
with $O(11^\rho+\rho)$ runs. Tape~$1$ crosses $\mathrm{U1}_1$ first; when it later reaches $\mathrm{LD}_1$, tapes $0$ and $2$ have wait age
far above $2$, so the taint of $\mathrm{LD}_1$ stays in $[db,3db]$. In the emit, the step that appends output $u$ is tainted only if
$B_0[u]$ or $\mathrm{LD}_1[u]$ is live tainted, if one of the tape-$2$ cells $3u,3u+1,3u+2$ below $E_2$ is live tainted, or if
$u=\lceil E_2/3\rceil$ and $B_0[u-1]$ is live tainted (the park of tape~$2$ can extend the tainted interval of $u-1$ by one output).
Within $[db,3db]$ the outputs of $\Out(A)$ are the $4b$ bits of the entries $(1,1)$, $(1,2)$, $(2,1)$ and $(2,2)$. Elsewhere, $B_0$ and the
tape-$2$ frames $Z_r$, $\mathrm{BS}_r$ of levels $r<\rho$ give $O(11^\rho)$ tainted intervals of outputs, of total length
$O(9^\rho db+18^\rho)$, each block of $db$ output positions contains at most $2b$ outputs of $\Out(A)$, and $(18/11)^\rho<2^\rho\le db$.
So
\begin{equation}\label{eq:c4-total}
  \abs{C_0}\le\frac{db}{2^\rho}+5,\qquad N_{\mathrm{run}}\le c''\Bigl(11^\rho+\frac{9^\rho}d+\frac{18^\rho}L+\rho\Bigr),\qquad
  N_{\mathrm{up}}\le4b+c'\,11^\rho b .
\end{equation}

\subsection{Proof of \texorpdfstring{\cref{thm:cascade-cuts}}{the theorem} and \texorpdfstring{\cref{rem:cascade-widths}}{the remark}}

Let $b=b(d)=k$, let $\nu$ be as in \cref{thm:cascade-cuts}, and choose $\rho$ with $\nu\le2^\rho<2\nu$; for $d\ge d_0$ this gives
$1\le\rho\le r_{\max}$ and $(5/3)^\rho\le db$. For Cascade~3, $6^\rho<6\nu^{\log_26}\le6(db)^{0.52}$,
$5^\rho/d\le5(db)^{0.47}/d$ and $40^\rho/L\le40(db)^{1.07}/(d^2b)$, and the last two tend to $0$. For Cascade~4,
$11^\rho<11\nu^{\log_211}\le11(db)^{0.58}$, $9^\rho/d\le9(db)^{0.53}/d$ and $18^\rho/L\le18(db)^{0.7}/(d^2b)$, and again the last two
tend to $0$. Since $b=\log_2d$, \eqref{eq:c3-total} and \eqref{eq:c4-total} give $\abs{C_0}$, $N_{\mathrm{run}}$, $N_{\mathrm{up}}\le c_1\,db/\nu$
for $d\ge d_0$, with an absolute constant $c_1$.

(a) By \cref{lem:closure}(b), $\Gamma$ separates $\Out(\Qstar)=\Out(A)$ from $\Lin(\{1,2\})$, and $\abs\Gamma\le3c_1\,db/\nu$. \Cref{lem:R}
for the oblivious machine ($\eta=0$) gives $(\wval+1)\abs\Gamma\ge\wval\abs\Gamma+\abs{\Gamma\cap\Out(\Qstar)}\ge2db$, so
$\wval+1\ge2\nu/(3c_1)$.

(b) Let $V$ be the set of $v\in A$ whose root $\Out(v)$ is untainted, so $\abs V\ge2db-N_{\mathrm{up}}\ge db$ for $d\ge d_0$. By
\cref{lem:closure}(a), every path from a read of $A$ to a step of $\Out(V)$ uses an edge of $C$, which has at most $\abs{C_0}+3N_{\mathrm{run}}$
cell edges and at most $N_{\mathrm{run}}$ state edges. \Cref{lem:cellcut} gives
$db\le4c_1(db/\nu)\bigl(\log_2(g+1)+\log_2(q+1)\bigr)$, so $\log_2(g+1)+\log_2(q+1)\ge\nu/(4c_1)$.

For \cref{rem:cascade-widths}, let $b\ge1$ be arbitrary and choose $d^{1/5}\le2^\rho<2d^{1/5}$ (Cascade~3) or $d^{1/6}\le2^\rho<2d^{1/6}$
(Cascade~4). The same estimates give $N_{\mathrm{run}}=O(d^{0.52})$, respectively $O(d^{0.58})$, and $N_{\mathrm{up}}\le4b+O(d^{0.52}b)$,
respectively $4b+O(d^{0.58}b)$, uniformly in $b$, while $\abs{C_0}\le d^{4/5}b+2$, respectively $d^{5/6}b+5$. So $\abs V\ge db$ for $d\ge d_0$,
and the proof of \cref{lem:cellcut}, which uses nothing about the entry width, excludes the trace once $\log_2(g+1)+\log_2(q+1)$ is below
$d^{1/5}/c$, respectively $d^{1/6}/c$.
